\documentclass[11pt]{article}
\usepackage{fontspec}
\usepackage[french]{babel}
\usepackage{amsmath}
\usepackage{amssymb}
\usepackage{amsthm}
\usepackage{mathrsfs}
\usepackage{tikz-cd}
\usepackage{enumitem}

\newcommand{\printedpage}[1]{\marginpar{\footnotesize p.~#1}}
\newcommand{\msp}[1]{\marginpar{\footnotesize [#1]}}
\newcommand{\uHom}{\underline{\mathrm{Hom}}}
\newcommand{\Fl}{\operatorname{Fl}}
\newcommand{\Ob}{\operatorname{Ob}}
\newcommand{\End}{\operatorname{End}}
\newcommand{\Ind}{\operatorname{Ind}}
\newcommand{\id}{\mathrm{id}}
\newcommand{\cardred}{\mathop{\textup{card\,réd}}}
\newcommand{\card}{\operatorname{card}}
\newcommand{\edbr}[1]{\textup{\texttt{#1}}}

\begin{document}

\printedpage{139}
% head: A. Grothendieck : Dérivateurs, Ch. XVIII, Édition des Universités de Montpellier II et Paris VII
% foot: Version 17 décembre 2025

\noindent\textbf{Corollaire 4.13.39.} {\itshape Sous les conditions de 4.13.38, supposons que $K$ contienne un objet $e$ maximal dans $J$. Soit $\Sigma$ l'ensemble des endomorphismes $p$ de $e$ dans $J$, tels que pour toute double-flèche $k \overset{u}{\underset{v}{\rightrightarrows}} e$ \textup{dans} $K$, on ait $pu=pv$. Alors $\Sigma$ est stable par composition, et il existe une partie non vide $\Sigma_\infty$ \edbr{[}$=\Sigma_0$\edbr{]} de $\Sigma$, telle que pour $p,p'\in\Sigma_0$, existe
\msp{page 367}
$q\in\Sigma_0$ tel que $qp=qp'$ (i.e. telle que la sous-catégorie $J_0$ de $J$ dont l'ensemble des objets est $\{e\}$, et l'ensemble des flèches $\{\id_e\}\cup\Sigma_\infty$, soit filtrante). Si $\Sigma_0$ est choisie ainsi, soit $\mathcal{F}$ l'ensemble réunion de $\Fl K$ et de l'ensemble des flèches composées
\[
k \xrightarrow{\ u\ } e \xrightarrow{\ p\ } e \xrightarrow{\ v\ } k'\ , \tag*{$(*)$}
\]
avec $u,v\in\Fl J$, $p\in\Sigma_\infty$. Alors $\mathcal{F}$ est l'ensemble des flèches d'une sous-catégorie $K'$ de $J$, ayant mêmes objets que $J$, et admettant $e$ pour plus grand élément. De plus, $K'$ contient $K$ et est filtrante.}

\noindent\textsc{Démonstration}. Que $\Sigma$ soit stable par composition est clair. Pour construire $\Sigma_\infty$ ayant les propriétés énoncées, on note que, $J$ étant grande devant $\pi$, $\Sigma$ est non vide (\footnote[162]{cf. argument plus bas pour construction de $p_i$.}). On construit par récurrence une suite d'éléments
\[
p_0\ ,\ p_1\ ,\ \ldots\ ,\ p_n\ ,\ \ldots
\]
de $\Sigma$, en prenant pour $p_0$ un élément quelconque de $\Sigma$, et pour $i\geqslant 1$, prenant $p_i\in\Sigma$ tel que pour toute double-flèche $(p,p')$ du sous-monoïde $M_{i-1}$ de $\End_J(e)$ engendré par $(p_j)_{0\leqslant j\leqslant i-1}$, on ait $p_ip=p_ip'$. Un tel $p_i$ existe, car prenant l'ensemble des double-flèches de la forme précédente $(p,p')$, avec $p,p'\in M_{i-1}$, et de celles de la forme $(u,v) : k \rightrightarrows e$, avec $u,v\in\Fl J$, l'ensemble de
\msp{page 368}
toutes ces double-flèches est de cardinal $\leqslant\pi$ ($M_{i-1}$ étant dénombrable, donc de cardinal $\leqslant\pi$), donc il existe une flèche $e\longrightarrow j$ dans $J$ qui les égalise toutes ($J$ étant grande devant $\pi$). Quitte à composer avec $j\longrightarrow e$ ($e$ étant objet maximal de $J$), on trouve le $p_i : e\longrightarrow e$ cherché.

Soit alors $\Sigma_0$ le sous-ensemble multiplicatif de $\Sigma$ engendré par les $p_i$ ($i\in\mathbf{N}$), limite inductive des ensembles similaires $\Sigma_i$ engendrés par les $p_j$ pour $j\in[0,i]$. Si $p$, $p'$ sont dans $\Sigma_\infty$, il sont dans un même $\Sigma_i$, donc on aura $p_{i+1}p=p_{i+1}p'$, ce qui prouve que $\Sigma_\infty$ a la propriété filtrante exigée. % [sic] il sont -> ils sont, agreement of the verb with the plural subject

Prouvons maintenant que $\mathcal{F}$ est l'ensemble de flèches d'une sous-catégorie $K'$ de $J$, ayant $\Ob K$ comme ensemble d'objets. Il reste à établir la stabilité par composition $vu$. Mais si $u$ et $v$ sont dans $\Fl K$, ou si l'un des deux est dans $\Fl K$, l'autre non (donc de la forme $(*)$), cette stabilité est évidente. Reste à voir le cas où ni \edbr{[}$u$\edbr{]} ni $v$ sont dans $\Fl K$, donc sont tous deux de la forme $(*)$ :
\[
k \xrightarrow{\ u\ } e \xrightarrow{\ p\ } e \xrightarrow{\ v\ } k' \xrightarrow{\ u'\ } e \xrightarrow{\ p'\ } e \xrightarrow{\ v'\ } k''\ ,
\]
avec $u,v,u',v'\in\Fl K$, $p,p'\in\Sigma_\infty$. Comme $(u'v,\id_e)$ est une double-flèche de
\printedpage{140}
% head: A. Grothendieck : Dérivateurs, Ch. XVIII, Édition des Universités de Montpellier II et Paris VII
% foot: Version 17 décembre 2025
\msp{page 369}
$K$ de but $e$, on a
\[
p'u'v \;=\; p'\id_e \;=\; p'\ ,
\]
donc le composé en question s'écrit aussi
\[
k \xrightarrow{\ u\ } e \xrightarrow{\ p'p\ } e \xrightarrow{\ v'\ } k''\ ,
\]
où $p'p\in\Sigma_\infty$ puisque $\Sigma_\infty$ est stable par composition, donc il est dans $\mathcal{F}$, q.e.d.

Il est clair que $K'$ contient $K$, et ayant mêmes objets que $K$ qui admet $e$ comme plus grand élément, que $e$ est un plus grand élément de $K'$. Reste à prouver que $K'$ est filtrante. Comme $K'$ a un plus grand élément, il reste à voir que toute double-flèche $k\rightrightarrows l$ dans $K'$ peut être égalisée dans $K'$. Quitte à composer avec $l\longrightarrow e$, on peut supposer $l=e$. Or on voit que toute flèche $k\longrightarrow e$ dans $K'$ est de la forme
\[
k \xrightarrow{\ u\ } e \xrightarrow{\ p\ } e \xrightarrow{\ v\ } e\ ,
\]
avec $u,v\in\Fl K$ et $p\in\Sigma_\infty\cup\{\id_e\}$. Quitte à composer avec $p_0\in\Sigma_\infty$ fixé, cette flèche devient (puisque $p_0v=p_0\id_e=p_0$)
\[
k \xrightarrow{\ u\ } e \xrightarrow{\ p_0p\ } e\ , \qquad\quad \text{i.e. } p_0pu\ .
\]
Pour une flèche $v'p'u' : k\longrightarrow e$, on trouve $p_0p'u'$. Soit alors $q\in\Sigma_0$ tel que $q(p_0p)=q(p_0p')$, soit $r$ ($q$ existe puisque $p_0p,p_0p'\in\Sigma_\infty$). Donc $(q(p_0pu),q(p_0p'u'))=(ru,ru')$, et on a $ru=ru'$ puisque $r\in\Sigma_\infty$. Au total, $qp_0$ égalise les deux flèches $vpu$ et $v'p'u'$, q.e.d.

\msp{page 370}

\section*{4.14\quad Critères d'accessibilité pour grosses catégories}

\noindent\textbf{Théorème 4.14.1.} {\itshape Soient $\mathcal{C}$ une catégorie équivalente à une petite catégorie, $\Phi$, $\Psi$ deux petits ensembles de catégories, équivalentes à des petites catégories. Pour toute catégorie $I$, on désigne par $I_{\bullet}$ (resp. ${}_{\bullet}I$) le cône droit (resp. gauche) sur $I$, déduit de $I$ en ajoutant un objet final (resp. initial) $e=e_I$. Soit $\Lambda_1$ (resp. $\Lambda_2$) un ensemble de couples $(I,\varphi)$ (resp. $(I,\psi)$), où $I\in\Phi$ (resp. $I\in\Psi$), et où
\[
\varphi\ :\ I \longrightarrow \mathcal{C} \qquad\quad \text{(resp. }\psi\ :\ I \longrightarrow \mathcal{C}\ \ \text{)} \tag*{\textup{(4.14.1.1)}}
\]
est un foncteur. Pour toute catégorie $\mathcal{N}$ stable (\footnote[163]{Cette propriété de stabilité n'est pas nécessaire pour la définition de $\uHom^{\Lambda_2}_{\Lambda_1}(\mathcal{C},\mathcal{N})$, et il est presque sûr que le théorème est valable sans faire cette hypothèse sur $\mathcal{N}$. Cf. 4.14.16, où on voit que c'est en tous cas inutile pour $\Psi$.}), par $\varinjlim$ de type $\Phi$ et par $\varprojlim$ de type $\Psi$ on désigne par
\[
\uHom^{\Lambda_2}_{\Lambda_1}(\mathcal{C},\mathcal{N})\ \subset\ \uHom(\mathcal{C},\mathcal{N}) \tag*{\textup{(4.14.1.2)}}
\]
la sous-catégorie strictement pleine de $\uHom(\mathcal{C},\mathcal{N})$ formée des foncteurs $f:\mathcal{C}\longrightarrow\mathcal{N}$ ayant la propriété suivante :
\begin{enumerate}[leftmargin=2.5em]
\item[\textit{a)}] $\forall\ (I,\varphi)\in\Lambda_1$, le foncteur composé
\[
I_{\bullet} \xrightarrow{\ \varphi\ } \mathcal{C} \xrightarrow{\ f\ } \mathcal{N}
\]
est $I$-exact à gauche, i.e. fait de $(f\varphi)(e_I)$ une $\varinjlim\limits_{I}$ de $f\varphi|I$.
\printedpage{141}
% head: A. Grothendieck : Dérivateurs, Ch. XVIII, Édition des Universités de Montpellier II et Paris VII
% foot: Version 17 décembre 2025
\msp{page 371}
\item[\textit{b)}] Pour tout $(I,\psi)\in\Lambda_2$, le foncteur composé
\[
{}_{\bullet}I \xrightarrow{\ \psi\ } \mathcal{C} \xrightarrow{\ f\ } \mathcal{N}
\]
est $I$-exact à droite, i.e. fait de $(f\psi)(e_I)$ une limite projective de $(f\psi)|I$.
\end{enumerate}
\noindent Avec ces notations : si $\mathcal{N}$ est une catégorie accessible, il en est de même de $\uHom^{\Lambda_2}_{\Lambda_1}(\mathcal{C},\mathcal{N})$.}

\noindent\textbf{Commentaire.} On utilisera ce théorème surtout quand on a $\Lambda_2=\emptyset$ ou $\Lambda_1=\emptyset$, auquel cas on note $\uHom_{\Lambda_1}(\mathcal{C},\mathcal{N})$, resp. $\uHom^{\Lambda_2}(\mathcal{C},\mathcal{N})$, la catégorie $\uHom^{\Lambda_2}_{\Lambda_1}(\mathcal{C},\mathcal{N})$. On appliquera le théorème surtout quand $\mathcal{C}$ elle-même est stable par $\varinjlim$ de type $\Phi$ et par $\varprojlim$ de type $\Psi$, et quand $\Lambda_1$ (resp. $\Lambda_2$) est déjà formé de foncteurs $\varphi$ (resp. $\psi$) qui sont $I$-exacts à gauche (resp. à droite). Donc $\uHom^{\Lambda_2}_{\Lambda_1}(\mathcal{C},\mathcal{N})$ est alors formée des $f:\mathcal{C}\longrightarrow\mathcal{N}$ qui commutent à certaines $\varinjlim$ de type $\Phi$, et à certaines $\varprojlim$ de type $\Psi$. En particulier, on trouve :

\msp{page 372}

\noindent\textbf{Corollaire 4.14.2.} {\itshape Soient $\mathcal{C}$ une catégorie équivalente à une petite catégorie, $\mathcal{N}$ une $\mathfrak{U}$-catégorie, $\Phi$ et $\Psi$ deux petits ensembles de catégories $I$, équivalentes à des petites catégories. On suppose $\mathcal{C}$ et $\mathcal{N}$ stables par $\varinjlim$ de type $\Phi$, par $\varprojlim$ de type $\Psi$. Soit $\uHom^{\Psi}_{\Phi}(\mathcal{C},\mathcal{N})$ la sous-catégorie strictement pleine de $\uHom(\mathcal{C},\mathcal{N})$ formée des $f:\mathcal{C}\longrightarrow\mathcal{N}$ qui commutent aux $\varinjlim$ de type $\Phi$ et aux $\varprojlim$ de type $\Psi$. Alors, si $\mathcal{N}$ est accessible, il en est de même de $\uHom^{\Psi}_{\Phi}(\mathcal{C},\mathcal{N})$.}

Voici des cas particuliers remarquables de ce corollaire :

\noindent\textbf{Corollaire 4.14.3.} {\itshape Soient $\mathcal{M}$, $\mathcal{N}$ deux catégories accessibles, et $\pi_0$ un cardinal infini, tel que $\mathcal{M}$ satisfasse à $L_{\pi_0}$. Soit $\uHom_{\pi_0}(\mathcal{M},\mathcal{N})$ la sous-catégorie strictement pleine de $\uHom(\mathcal{M},\mathcal{N})$ formée des foncteurs $\pi_0$-accessibles. Alors $\uHom_{\pi_0}(\mathcal{M},\mathcal{N})$ est accessible.}

\msp{page 373}

(Le cas où $\pi_0$ est utile pour $\mathcal{M}$ a déjà été traité.)

\noindent\textsc{Démonstration}. Dans le cas général, soit $\pi>\pi_0$ un cardinal utile pour $\mathcal{M}$, de sorte qu'on a
\[
\mathcal{M} \xleftarrow{\ \sim\ } \Ind_\pi(\mathcal{C}) \qquad\quad \text{avec } \mathcal{C}=\mathcal{M}_\pi\ .
\]
On peut supposer $\mathcal{M}=\Ind_\pi(\mathcal{C})$. Le fait que $\mathcal{M}$ satisfasse à $\mathrm{L}_{\pi_0}$ implique que $\mathcal{C}$ est stable par $\varinjlim$ grandes devant $\pi_0$, et de cardinal $\leqslant\pi$. Soit $\Phi$ l'ensemble des catégories de cette nature. $\Phi$ n'est pas petit, mais il existe $\Phi_0\subset\Phi$ petit, tel que toute $I\in\Phi$ soit isomorphe à une $I_0\in\Phi_0$, ce qui suffira pour la validité de 4.14.1 et 4.14.2 (où on fera ici $\Psi=\emptyset$). On peut considérer $\uHom_{!\,\pi_0}(\mathcal{M},\mathcal{N})$ comme une sous-catégorie strictement pleine de $\uHom_{!\,\pi}(\mathcal{M},\mathcal{N})\xrightarrow{\ \sim\ }\uHom(\mathcal{C},\mathcal{N})$ (cf. 4.12.3, p. 275). Pour un foncteur $f:\mathcal{M}\longrightarrow\mathcal{N}$ qui est $\pi$-accessible, dire qu'il est même $\pi_0$-accessible revient à dire que sa restriction à $\mathcal{C}$, i.e. son image dans $\uHom(\mathcal{C},\mathcal{N})$, commute aux $\varinjlim$ de type $\Phi$. (Cf. 4.13.20, page 341 (\footnote[164]{La référence n'est pas tout à fait adéquate. Il faudrait réunir dans un paragraphe à part tous les énoncés de ce type, dont certains utilisent le théorème de représentation indicielle.}), où on prend $\alpha=\alpha_\pi$, notant que $J_{\alpha_\pi}$ grande devant $\pi_0$ puisque $\pi>\pi_0$).

\msp{page 374}

Donc on trouve
\[
\uHom_{!\,\pi_0}(\mathcal{M},\mathcal{N}) \xrightarrow{\ \sim\ } \uHom_{\Phi}(\mathcal{C},\mathcal{N})\ ,
\]
\printedpage{142}
% head: A. Grothendieck : Dérivateurs, Ch. XVIII, Édition des Universités de Montpellier II et Paris VII
% foot: Version 17 décembre 2025
et le corollaire 4.14.2 donne la conclusion voulue.

\noindent\textbf{Corollaire 4.14.4.} {\itshape Soient $\mathcal{M}$, $\mathcal{N}$ deux catégories accessibles, $\Phi$ et $\Psi$ deux ensembles des catégories comme dans 4.14.2, on suppose que $\mathcal{M}$ et $\mathcal{N}$ satisfassent $L_{\pi_0}$, et \edbr{[soient]} stables par $\varinjlim$ de type $\Phi$ et par $\varprojlim$ de type $\Psi$. Soit $\uHom^{\Psi}_{\Phi,\pi_0}(\mathcal{M},\mathcal{N})$ la sous-catégorie strictement pleine de $\uHom(\mathcal{M},\mathcal{N})$ formée des foncteurs $\pi_0$-accessibles qui commutent aux $\varinjlim$ de type $\Phi$, et aux $\varprojlim$ de type $\Psi$. Alors $\uHom^{\Psi}_{\Phi,\pi_0}(\mathcal{M},\mathcal{N})$ est accessible.}

\noindent\textsc{Démonstration}. Cette fois, on choisit $\pi>\pi_0$ tel qu'on ait $\cardred I\leqslant\pi$ pour tout $I$ dans $\Phi$ ou dans $\Psi$. La catégorie envisagée s'interprète alors (posant encore $\mathcal{C}=\mathcal{M}_\pi$) comme
\msp{page 375}
\[
\uHom^{\Psi}_{\Phi,\pi_0}(\mathcal{M},\mathcal{N}) \xrightarrow{\ \sim\ } \uHom^{\Psi}_{\Phi'}(\mathcal{C},\mathcal{N})\ ,
\]
où $\Phi'$ est formée des catégories qui sont soit dans $\Phi$, ou qui sont grandes devant $\pi_0$ et de cardinal $\leqslant\pi$. Il suffit donc d'appliquer encore 4.14.2.

\noindent\textbf{Corollaire 4.14.5.} {\itshape Considérons les trois ensembles suivants :
\begin{enumerate}[leftmargin=2.5em]
\item[\textit{a)}] $\Phi_1$ est formée de toutes les catégories équivalentes à des petites catégories.
\item[\textit{b)}] $\Phi_2$ comme a), mais en restreignant aux catégories filtrantes.
\item[\textit{c)}] $\Phi_3$ formé des catégories discrètes, et de celles qui sont grandes devant $\pi_0$ ($\pi_0$ un cardinal infini donné).
\end{enumerate}
\noindent Soit $\Phi$ l'un des ensembles de catégories $\Phi_1$, $\Phi_2$, $\Phi_3$, $\Phi_2\cup\Phi_3$ ou la réunion d'un de ceux-ci et d'un petit ensemble $\Phi_0$ de catégories essentiellement petites (\footnote[165]{On pourrait faire ici un énoncé qui englobe 4.14.4 en rajoutant un
\begin{enumerate}[leftmargin=2.5em]
\item[d)] $\Phi_4$ formé des catégories grandes devant $\pi_0$,
\end{enumerate}
\noindent et prendre $\Phi$ de la forme de $\Phi_1$, $\Phi_2$, $\Phi_4$, $\Phi_3$, $\Phi_3\cup\Phi_2$, $\Phi_3\cup\Phi_4$, ou la réunion de ceux-ci et d'un $\Phi_0$.}). Soit $\Psi$ un petit ensemble de catégories. Soient $\mathcal{M}$, $\mathcal{N}$ des catégories accessibles, stables par $\varinjlim$ de type $\Phi$, et par $\varprojlim$ de type $\Psi$. Alors $\uHom^{\Psi}_{\Phi}(\mathcal{M},\mathcal{N})$ est accessible.}

\msp{page 376}

\noindent\textsc{Démonstration}. Soit $\pi'_0\geqslant\pi_0$ un cardinal utile pour $\mathcal{M}$, $\pi$ un cardinal ayant les propriétés suivantes.
\begin{enumerate}[leftmargin=2.5em]
\item[a)] $\pi\geqslant\sup(\pi'_0,\cardred\mathcal{M}_{\pi'_0})$, $\pi>\pi_0$.
\item[b)] $\pi^{\pi'_0}=\pi$.
\item[c)] $I\in\Phi_0\cup\Psi\Longrightarrow\cardred I\leqslant\pi$.
\end{enumerate}
\noindent Soit $\mathcal{C}=\mathcal{M}_\pi$, d'où $\mathcal{M}\simeq\Ind_\pi(\mathcal{C})$. Dans tous les cas envisagés, $\Phi$ contient l'ensemble des catégories essentiellement petites qui sont grandes devant $\pi_0$, a fortiori celles qui sont grandes devant $\pi$, donc la catégorie $\uHom^{\Psi}_{\Phi}$ s'interprète comme une sous-catégorie strictement pleine de
\[
\uHom(\mathcal{C},\mathcal{N})\ .
\]
Si d'autre part $\Phi_\pi$ désigne l'ensemble des $I\in\Phi$ telles que $\card I\leqslant\pi$, on trouve
\[
\uHom^{\Psi}_{\Phi}(\mathcal{M},\mathcal{N})\ \simeq\ \uHom^{\Psi}_{\Phi_\pi}(\mathcal{C},\mathcal{N})\ ,
\]
et on applique 4.14.2.

\printedpage{143}
% head: A. Grothendieck : Dérivateurs, Ch. XVIII, Édition des Universités de Montpellier II et Paris VII
% foot: Version 17 décembre 2025
\msp{page 377}

En particulier, on trouve :

\noindent\textbf{Corollaire 4.14.6.} {\itshape Soient $\mathcal{M}$, $\mathcal{N}$ deux catégories accessibles, supposons $\mathcal{M}$ stable par petites $\varinjlim$ (i.e. $\mathcal{M}$ est une catégorie paratopos). Alors la catégorie $\uHom_{!}(\mathcal{M},\mathcal{N})$ des foncteurs $f:\mathcal{M}\longrightarrow\mathcal{N}$ qui commutent aux petites $\varinjlim$ est accessible. Si $\mathcal{N}$ est elle aussi un paratopos, $\uHom_{!}(\mathcal{M},\mathcal{N})$ est un paratopos.}

\noindent\textsc{Démonstration}. La première assertion est le cas particulier de 4.14.5 relatif au cas où $\Phi=\Phi_0$, $\Psi=\emptyset$. La deuxième assertion en résulte, puisque $\uHom_{!}(\mathcal{M},\mathcal{N})$ est visiblement stable par petites $\varinjlim$ , donc c'est un paratopos étant accessible (4.8.9, page 207).

\noindent\textbf{Corollaire 4.14.7.} {\itshape Soient $\Phi$, $\Psi$, $\mathcal{M}$ comme dans 4.14.5. Alors pour toute catégorie accessible $\mathcal{N}$, la catégorie
\msp{page 378}
\[
\uHom^{\Phi}_{\Psi}(\mathcal{M}^{\circ},\mathcal{N})\ \simeq\ (\uHom^{\Psi}_{\Phi}(\mathcal{M},\mathcal{N}^{\circ}))^{\circ}
\]
est accessible.}

\noindent\textsc{Démonstration}. On procède comme pour 4.14.5, pour expliciter
\[
\uHom^{\Psi}_{\Phi}(\mathcal{M},\mathcal{N}^{\circ}) \xrightarrow{\ \sim\ } \uHom^{\Psi}_{\Phi_\pi}(\mathcal{C},\mathcal{N}^{\circ})\ ,
\]
d'où
\[
\uHom^{\Phi}_{\Psi}(\mathcal{M}^{\circ},\mathcal{N}) \xrightarrow{\ \sim\ } \uHom^{\Phi_\pi}_{\Psi}(\mathcal{C}^{\circ},\mathcal{N})\ ,
\]
et on applique 4.14.2 à $\mathcal{C}^{\circ}$, $\Psi$, $\Phi_\pi$.

\noindent\textbf{Corollaire 4.14.8.} {\itshape Soient $\mathcal{M}$ et $\mathcal{N}$ deux catégories paratopos. Alors
\[
\mathcal{M}\boxtimes\mathcal{N}\ \simeq\ \uHom^{!}(\mathcal{M}^{\circ},\mathcal{N}) \qquad\quad \text{(cf. 4.2 et 4.3)}
\]
est un paratopos.}

En effet, en vertu de 4.14.7 (avec $\Phi=\Phi_\pi$ [?], $\Psi=\emptyset$), c'est une catégorie accessible.

\msp{page 379}

D'autre part, il est immédiat qu'elle est stable par petites $\varprojlim$ , $\mathcal{N}$ l'étant. Donc on conclut encore par 4.8.9 (page 207) que c'est une catégorie paratopos.

\noindent\textbf{Corollaire 4.14.9.} {\itshape Soient $\mathcal{M}$, $\mathcal{N}$ deux catégories paratopos, et $\uHom^{!}(\mathcal{M},\mathcal{N})$ la catégorie des foncteurs $\mathcal{M}\longrightarrow\mathcal{N}$ qui commutent aux petites $\varprojlim$ . On a alors une équivalence de catégories
\[
\begin{aligned}
\uHom^{!}(\mathcal{M},\mathcal{N}) &\ \xrightarrow{\ \sim\ }\ \uHom_{!}(\mathcal{N},\mathcal{M})^{\circ}\\
f &\ \longmapsto\ \text{adjoint à gauche de } f\ ,
\end{aligned}
\]
\noindent\textup{\textbf{!}}\ \textup{\textbf{!}} et par suite la catégorie \textup{\textbf{opposée}} à $\uHom^{!}(\mathcal{M},\mathcal{N})$ est accessible.} % Verified against PDF leaf 25: a nonverbal marginal curved mark; see the correction register.

Ceci résulte du

\noindent\textbf{Lemme 4.14.10} (pour mémoire). {\itshape Soient $\mathcal{M}$ \edbr{[une]} catégorie pseudo-topos, $f:\mathcal{M}\longrightarrow\mathcal{N}$ un foncteur. Pour que $f$ admette un adjoint à gauche (resp. à droite), il faut et il suffit que $f$ commute aux petites $\varprojlim$ (resp. aux petites $\varinjlim$ ).}

\end{document}
